\documentclass[a4paper,11pt]{article}
\usepackage[a4paper,margin=2cm]{geometry}
\usepackage[T1]{fontenc}
\usepackage[utf8]{inputenc}
\usepackage[ngerman,english]{babel}
\usepackage{lmodern,microtype,xcolor,parskip,enumitem}
\usepackage[most]{tcolorbox}
\usepackage{titlesec}
\usepackage[hidelinks]{hyperref}
\definecolor{HeaderBlueDark}{RGB}{70,110,170}
\definecolor{RowBlueLight}{RGB}{235,243,255}
\definecolor{RowBlueDark}{RGB}{220,233,250}
\definecolor{SoftGray}{RGB}{90,90,90}
\setlength{\parindent}{0pt}
\setlist[itemize]{leftmargin=1.4em,itemsep=0.35em,topsep=0.35em}
\linespread{1.06}
\titleformat{\section}[block]{\Large\bfseries\color{HeaderBlueDark}}{}{0pt}{}
\titlespacing{\section}{0pt}{1.3em}{0.8em}
\newtcolorbox{exbox}{enhanced,breakable,colback=RowBlueLight,colframe=RowBlueDark,boxrule=0.4pt,arc=1.5mm,left=1.2mm,right=1.2mm,top=1mm,bottom=1mm,title={Beispiele \textnormal{(Examples)}},fonttitle=\bfseries,colbacktitle=RowBlueDark,coltitle=black}
\NewDocumentEnvironment{grammarentry}{mm+b}{\begin{tcolorbox}[enhanced,breakable,colback=white,colframe=HeaderBlueDark,boxrule=0.6pt,arc=1.5mm,left=1.5mm,right=1.5mm,top=1mm,bottom=1mm,colbacktitle=HeaderBlueDark,coltitle=white,fonttitle=\bfseries,title={#1\hfill\normalfont\footnotesize #2}]#3\end{tcolorbox}}{}
\newcommand{\GermanText}[1]{\textbf{Deutsch: }#1\par}
\newcommand{\EnglishText}[1]{{\small\color{SoftGray}\textbf{English: }#1\par}}
\newcommand{\ExampleItem}[2]{\item \textbf{#1}\\{\small\color{SoftGray}#2}}
\title{Es gibt}
\date{}
\begin{document}
   \maketitle
   \tableofcontents
   \newpage
   \section{Bedeutung (Meaning)}
      \begin{grammarentry}{es gibt}{Existenzkonstruktion / existential construction}
         \GermanText{Die Verbindung \textbf{es gibt} bezeichnet Existenz, Vorhandensein oder Verfügbarkeit. Das \textbf{es} ist ein formales Subjekt, kein Verweis auf eine bestimmte Sache.}
         \EnglishText{The expression means there is / there are. The pronoun es is a formal subject, not a reference to a specific object.}
         \GermanText{Grundmuster: \textbf{Es gibt + Akkusativ}. Das Verb bleibt im Singular, auch wenn mehrere Dinge vorhanden sind.}
         \EnglishText{Pattern: es gibt + accusative. The verb remains singular even if the object is plural.}
      \end{grammarentry}
      \begin{exbox}\begin{itemize}
         \ExampleItem{Es gibt einen Bahnhof.}{There is a station.}
         \ExampleItem{Es gibt eine Haltestelle.}{There is a stop.}
         \ExampleItem{Es gibt ein Museum.}{There is a museum.}
         \ExampleItem{Es gibt viele Museen.}{There are many museums.}
         \end{itemize}\end{exbox}
   \section{Kasus und Rektion (Case and government)}
         \begin{grammarentry}{Akkusativ, nicht Dativ}{accusative, not dative}
            \GermanText{Das Existierende ist das \textbf{Akkusativobjekt}: \textbf{einen} Mann, \textbf{eine} Frau, \textbf{ein} Kind, \textbf{viele} Menschen.}
            \EnglishText{The thing whose existence is stated takes the accusative case.}
            \GermanText{Ortsangaben haben ihren eigenen Kasus: \textbf{in der Nähe} steht im Dativ, weil \textbf{in} hier einen Ort beschreibt. Dieser Dativ kommt nicht von \textbf{es gibt}.}
            \EnglishText{Location phrases have a case of their own: in der Nähe is dative due to locative in, not due to es gibt.}
         \end{grammarentry}
         \begin{exbox}\begin{itemize}
            \ExampleItem{Es gibt in der Nähe einen Parkplatz.}{There is a parking space nearby.}
            \ExampleItem{Es gibt auf dem Tisch ein Buch.}{There is a book on the table.}
            \ExampleItem{In unserer Stadt gibt es eine Universität.}{There is a university in our city.}
            \end{itemize}\end{exbox}
   \section{Wortstellung (Word order)}
            \begin{grammarentry}{Satzbau}{sentence structure}
               \GermanText{Im Hauptsatz steht das konjugierte Verb auf Position zwei: \textbf{Hier gibt es ...}. Bei Ja/Nein-Fragen steht es vorne: \textbf{Gibt es ...?} Im Nebensatz mit \textbf{dass} steht es am Ende: \textbf{..., dass es ... gibt}.}
               \EnglishText{Main clauses are verb-second; yes/no questions verb-first; clauses introduced by dass are verb-final.}
            \end{grammarentry}
            \begin{exbox}\begin{itemize}
               \ExampleItem{Es gibt hier eine Schule.}{There is a school here.}
               \ExampleItem{Hier gibt es eine Schule.}{There is a school here.}
               \ExampleItem{Gibt es hier eine Schule?}{Is there a school here?}
               \ExampleItem{Ich weiß, dass es hier eine Schule gibt.}{I know that there is a school here.}
               \end{itemize}\end{exbox}
   \section{Zeitformen (Tenses)}
               \begin{grammarentry}{Formen}{conjugation}
                  \GermanText{Präsens: \textbf{es gibt}; Präteritum: \textbf{es gab}; Perfekt: \textbf{es hat ... gegeben}; Plusquamperfekt: \textbf{es hatte ... gegeben}; Futur I: \textbf{es wird ... geben}; Modalverb: \textbf{es kann ... geben}.}
                  \EnglishText{Present, simple past, present perfect, past perfect, future, and modal constructions.}
               \end{grammarentry}
               \begin{exbox}\begin{itemize}
                  \ExampleItem{Früher gab es hier ein Kino.}{There used to be a cinema here.}
                  \ExampleItem{Es hat einen Unfall gegeben.}{There has been an accident.}
                  \ExampleItem{Es hatte mehrere Warnungen gegeben.}{There had been several warnings.}
                  \ExampleItem{Morgen wird es ein Treffen geben.}{There will be a meeting tomorrow.}
                  \ExampleItem{Es könnte Probleme geben.}{There could be problems.}
                  \end{itemize}\end{exbox}
   \section{Verneinung und Mengen (Negation and quantities)}
                  \begin{grammarentry}{kein / nicht}{negation}
                     \GermanText{Vor einem verneinten Substantiv benutzt man meistens \textbf{kein} im Akkusativ: \textbf{keinen Parkplatz}, \textbf{keine Haltestelle}, \textbf{kein Museum}. Mit \textbf{nicht} kann man andere Teile verneinen.}
                     \EnglishText{Use accusative forms of kein to deny the existence of something. Nicht can negate other parts of a sentence.}
                  \end{grammarentry}
                  \begin{exbox}\begin{itemize}
                     \ExampleItem{Es gibt keinen freien Platz.}{There is no free space.}
                     \ExampleItem{Es gibt keine Busse.}{There are no buses.}
                     \ExampleItem{Es gibt nur wenige Möglichkeiten.}{There are only a few possibilities.}
                     \ExampleItem{Nicht überall gibt es gute Verbindungen.}{Not everywhere has good transport connections.}
                     \end{itemize}\end{exbox}
   \section{Vergleich mit sein (Comparison with sein)}
                     \begin{grammarentry}{Existenz und Standort}{existence and location}
                        \GermanText{\textbf{Es gibt} führt etwas als existierend ein. \textbf{Sein} gibt häufig den Ort oder Zustand von etwas Bekanntem an. Nach \textbf{sein} steht kein Akkusativobjekt.}
                        \EnglishText{Es gibt often introduces existence; sein usually locates or describes something already known. Sein does not take an accusative object.}
                     \end{grammarentry}
                     \begin{exbox}\begin{itemize}
                        \ExampleItem{In Graz gibt es ein Museum.}{There is a museum in Graz.}
                        \ExampleItem{Das Museum ist im Zentrum.}{The museum is in the centre.}
                        \end{itemize}\end{exbox}
   \section{Typische Wendungen (Common expressions)}
                        \begin{exbox}\begin{itemize}
                           \ExampleItem{Es gibt keinen Grund zur Sorge.}{There is no reason to worry.}
                           \ExampleItem{Es gibt nichts zu essen.}{There is nothing to eat.}
                           \ExampleItem{Es gibt heute Suppe.}{There is soup today.}
                           \ExampleItem{Was gibt's Neues?}{What's new?}
                           \end{itemize}\end{exbox}
   \section{Häufige Fehler (Common mistakes)}
                           \begin{exbox}\begin{itemize}
                              \ExampleItem{Falsch: Es gibt einem Parkplatz. Richtig: Es gibt einen Parkplatz.}{Accusative is required, not dative.}
                              \ExampleItem{Falsch: Hier es gibt einen Park. Richtig: Hier gibt es einen Park.}{The finite verb is second in a main clause.}
                              \ExampleItem{Falsch: ..., dass es gibt einen Park. Richtig: ..., dass es einen Park gibt.}{The finite verb is final in a dass-clause.}
                              \ExampleItem{Falsch: Es geben drei Schulen. Richtig: Es gibt drei Schulen.}{The verb stays singular.}
                              \end{itemize}\end{exbox}
   \section{Zusammenfassung (Summary)}
                              \GermanText{Merksatz: \textbf{es gibt + Akkusativ}, nicht Dativ. Ortsangaben haben einen eigenen Kasus. \textbf{Hier gibt es ...}; \textbf{..., dass es ... gibt}; \textbf{es gab}; \textbf{es gibt keinen/keine/kein ...}.}
                              \EnglishText{Remember: accusative object; location phrases follow their own case rules; verb-second in main clauses, verb-final in dass-clauses; past es gab; negation with kein.}
\end{document}
